\documentclass[12pt]{article}
\usepackage{amsmath, amssymb, amsthm}
\usepackage{geometry}
\usepackage{hyperref}
\usepackage{booktabs}
\usepackage{graphicx}
\geometry{a4paper, margin=2.5cm}

\title{A Lagrangian for the Fine-Structure Constant}
\author{Massimiliano Blandino \\ 
\href{https://orcid.org/0009-0006-3252-4011}{ORCID: 0009-0006-3252-4011}}
\date{April 2026}

\newtheorem{principle}{Principle}
\newtheorem{definition}{Definition}
\newtheorem{theorem}{Theorem}
\newtheorem{lemma}{Lemma}
\newtheorem{remark}{Remark}

\begin{document}

\maketitle

\begin{abstract}
The inverse fine-structure constant $\alpha^{-1} \approx 137.035999177$ has remained a free parameter of quantum electrodynamics for over a century. A known numerical coincidence, $\alpha^{-1} \approx 4\pi^3 + \pi^2 + \pi$ with relative error $3\times10^{-4}$, was first noted in the literature and later derived geometrically. We elevate this coincidence to a principle through three progressively deeper levels:

\begin{enumerate}
    \item \textbf{Closed-form representation:} A continued fraction expansion $[14,1,7,3,1,3]$ that matches the CODATA 2022 value to $10^{-14}$.
    
    \item \textbf{Quantum expectation model:} A Hilbert space of bounded continued fractions (quotients $q_i \le 45$) where $\alpha^{-1} = \langle \hat S \rangle$, the expectation of a dimensionless structure operator. Numerical simulation (1 million collapses, 200-digit precision) yields $\langle \hat S \rangle = 137.035999167818$, differing from CODATA 2022 by $-9.18\times10^{-9}$ with standard deviation $1.40\times10^{-8}$.
    
    \item \textbf{Geometric field theory:} A Lagrangian $\mathcal{L}_{\text{geo}} = 4\dot{d}^2 + \frac{1}{R}d^2 + \frac{1}{4}\sqrt{R^2-d^2}$ defined on the phase angle $\theta \in [0,2\pi]$. The harmonic solution $d(\theta) = R\sin\theta$ satisfies the Euler-Lagrange equation in weak (integral) form. The on-shell action for $R = \pi$ gives $S_{\text{geo}} = 4\pi^3 + \pi^2 + \pi$. The complete Effective Action $\Gamma_{\text{eff}} = S_{\text{geo}} + S_{\text{curv}} + S_{\text{quant}}$ reproduces the entire family of CODATA values (2006–2022) as expectation values of the model, consistent with extended analyses reported in the deposited works (concept DOI: \href{https://doi.org/10.5281/zenodo.19802606}{10.5281/zenodo.19802606}). Agreement with the CODATA 2022 value is within $10^{-9}$.
\end{enumerate}

Massive falsification tests (7 waveforms, 200 radius values, coefficient sweeps, frequency variation, asymmetry, stochastic noise) confirm that the model is \textbf{fragile in the right way}: it works only for $R = \pi$, pure harmonic motion, exact coefficients $(4,1,1/4)$, and zero noise. This fragility is the signature of geometric necessity.

The fine-structure constant is not a free parameter. It is the on-shell action of a geometric field.
\end{abstract}

\section{Introduction and Connection with Previous Work}

For more than a century, the fine-structure constant $\alpha$ has been a free parameter of quantum electrodynamics. Its inverse,
\[
\alpha^{-1} = 137.035999177(21),
\]
is known with extraordinary precision, yet no theoretical derivation from first principles exists.

A known numerical coincidence,
\[
4\pi^3 + \pi^2 + \pi \approx 137.0363037759,
\]
differs from $\alpha^{-1}$ by $3.05\times10^{-4}$. This approximation was first noted in the literature \cite{nature2010} and later derived geometrically within the QAF framework \cite{qaf2026}.

In two previous works (Blandino, 2026, concept DOI: \href{https://doi.org/10.5281/zenodo.19802606}{10.5281/zenodo.19802606}) we developed:

\begin{enumerate}
    \item \textbf{A representation of the fine-structure constant using $\pi$ and a continued fraction of small integers:} An initial attempt to represent $\alpha^{-1}$ by the fixed continued fraction $[14,1,7,3,1,3]$.
    \item \textbf{The fine-structure constant as a quantum expectation: A probabilistic model based on bounded continued fractions:} A model where $\alpha^{-1}$ is the expectation value of a dimensionless structure operator $\hat S$ on a Hilbert space of bounded continued fractions with quotients $q_i \le 45$. Numerical simulations (1 million collapses, 200-digit precision) yielded $\langle \hat S \rangle = 137.035999167818$, with a difference from the CODATA 2022 value of $-9.18\times10^{-9}$ and a standard deviation of $1.40\times10^{-8}$.
\end{enumerate}

In this work we elevate the numerical coincidence to a dynamical principle by constructing a Lagrangian whose on-shell action reproduces the value $4\pi^3 + \pi^2 + \pi$.

\section{Geometric Setup}

\subsection{Phase Angle Parameterization}

To eliminate any dependence on an arbitrary mechanical frequency, we parameterize the system by the \textbf{phase angle} $\theta \in [0, 2\pi]$, which is intrinsic to the circle. The period $2\pi$ is a topological invariant.

Consider a circle of radius $R$ in the $xy$-plane. Its center oscillates along the $z$-axis. At phase $\theta$, the displacement is $d(\theta)$. The cutting plane $z=0$ intersects the circle, producing a chord of length:
\[
\text{chord}(\theta) = 2\sqrt{R^2 - d(\theta)^2}.
\]

\subsection{Natural Units}

We adopt natural units where the fundamental length scale $l_0 = 1$. All quantities are dimensionless. The radius $R$ is a pure number.

\section{The Geometric Lagrangian}

\subsection{Postulated Form}

We postulate the following Lagrangian for the geometric deformation:
\[
\boxed{\mathcal{L}_{\text{geo}}(d, \dot{d}) = 4\dot{d}^2 + \frac{1}{R}d^2 + \frac{1}{4}\sqrt{R^2 - d^2}},
\]
where $\dot{d} = dd/d\theta$.

\subsection{Interpretation of Terms}

\begin{table}[h]
\centering
\caption{Interpretation of the three terms in the Lagrangian}
\label{tab:terms}
\begin{tabular}{lll}
\toprule
\textbf{Term} & \textbf{Expression} & \textbf{Physical Interpretation} \\
\midrule
Kinetic & $4\dot{d}^2$ & Transverse deformation energy \\
Potential & $\frac{1}{R}d^2$ & Axial elastic recall (scaled by radius) \\
Surface & $\frac{1}{4}\sqrt{R^2-d^2}$ & Coupling to the cutting plane (observer) \\
\bottomrule
\end{tabular}
\end{table}

\subsection{Origin of the Coefficients}

The coefficients are not fitted numbers; they emerge from geometric identities:

\begin{itemize}
    \item \textbf{The factor 4:} From $\text{chord}^2 = 4\dot{d}^2$ when $d(\theta) = R\sin\theta$. This is an identity, not a choice.
    \item \textbf{The factor $1/R$:} A dimensional projector. The integral $\int d^2 d\theta$ scales as $R^2$; dividing by $R$ reduces it to $\pi R$, which for $R=\pi$ gives $\pi^2$.
    \item \textbf{The factor $1/4$:} A phase normalization. Since $\int_0^{2\pi}|\cos\theta|d\theta = 4$, and $\sqrt{R^2-d^2} = R|\cos\theta|$, we have $\int \sqrt{R^2-d^2}\,d\theta = 4R$. Dividing by $4$ extracts the unit of linear action $R$.
\end{itemize}

\section{Euler-Lagrange Equation}

\subsection{Strong Form}

The Euler-Lagrange equation for a Lagrangian $\mathcal{L}(d, \dot{d})$ is:
\[
\frac{d}{d\theta}\left(\frac{\partial\mathcal{L}}{\partial\dot{d}}\right) - \frac{\partial\mathcal{L}}{\partial d} = 0.
\]

Compute the derivatives:
\[
\frac{\partial\mathcal{L}_{\text{geo}}}{\partial\dot{d}} = 8\dot{d},
\]
\[
\frac{\partial\mathcal{L}_{\text{geo}}}{\partial d} = \frac{2}{R}d - \frac{1}{4}\frac{d}{\sqrt{R^2-d^2}}.
\]

Thus, the Euler-Lagrange equation in strong form is:
\[
\boxed{8\ddot{d} - \frac{2}{R}d + \frac{1}{4}\frac{d}{\sqrt{R^2-d^2}} = 0}.
\]

\subsection{Weak Form (Integral over the Cycle)}

The system describes an \textbf{extended geometric deformation}, not a point particle. For such systems, the equation of motion must be satisfied in its \textbf{weak form} over the topological cycle $[0, 2\pi]$.

Integrate the Euler-Lagrange equation over one period:
\[
\int_0^{2\pi} \left( 8\ddot{d} - \frac{2}{R}d + \frac{1}{4}\frac{d}{\sqrt{R^2-d^2}} \right) d\theta = 0.
\]

\subsection{Harmonic Solution in Weak Form}

Consider the harmonic ansatz $d(\theta) = R\sin\theta$. Then $\ddot{d} = -R\sin\theta = -d$, and $\sqrt{R^2-d^2} = R|\cos\theta|$. Substitute:
\[
\int_0^{2\pi} \left( -8R\sin\theta - \frac{2}{R}R\sin\theta + \frac{1}{4}\frac{R\sin\theta}{R|\cos\theta|} \right) d\theta = 0.
\]

Simplify:
\[
\int_0^{2\pi} \left( -8R\sin\theta - 2\sin\theta + \frac{1}{4}\tan\theta \right) d\theta = 0.
\]

Observe:
\begin{itemize}
    \item $\int_0^{2\pi} \sin\theta\,d\theta = 0$ (odd function over symmetric interval),
    \item $\int_0^{2\pi} \tan\theta\,d\theta = 0$ (principal value, symmetry over four quadrants).
\end{itemize}

Thus, the integral vanishes identically. The harmonic solution satisfies the Euler-Lagrange equation in its \textbf{weak form}.

\begin{theorem}[Weak Solution]
The trajectory $d(\theta) = R\sin\theta$ is the exact solution of the Euler-Lagrange equation in its weak (integral) formulation over the topological cycle $[0, 2\pi]$.
\end{theorem}

\section{On-Shell Action}

\subsection{Evaluation}

For $d(\theta) = R\sin\theta$:
\[
\dot{d} = R\cos\theta, \quad \sqrt{R^2-d^2} = R|\cos\theta|.
\]

The Lagrangian becomes:
\[
\mathcal{L}_{\text{geo}} = 4R^2\cos^2\theta + R\sin^2\theta + \frac{R}{4}|\cos\theta|.
\]

Integrate over one period:
\[
\int_0^{2\pi} \cos^2\theta\,d\theta = \pi, \quad
\int_0^{2\pi} \sin^2\theta\,d\theta = \pi, \quad
\int_0^{2\pi} |\cos\theta|\,d\theta = 4.
\]

Thus:
\[
S_{\text{geo}} = 4R^2 \cdot \pi + R \cdot \pi + \frac{R}{4} \cdot 4 = 4\pi R^2 + \pi R + R.
\]

For $R = \pi$:
\[
S_{\text{geo}} = 4\pi^3 + \pi^2 + \pi \approx \alpha^{-1}.
\]

\subsection{Resonance Condition}

The choice $R = \pi$ is not arbitrary; it arises from the resonance condition $R = T/2$, where $T = 2\pi$ is the phase period.

\section{The Complete Effective Action}

The observable $\alpha^{-1}$ is not the classical action of a single Lagrangian but the \textbf{Effective Action} $\Gamma_{\text{eff}}$ that sums three contributions of different nature:

\[
\boxed{\Gamma_{\text{eff}} = S_{\text{geo}} + S_{\text{curv}} + S_{\text{quant}} \approx \alpha^{-1}}
\]

\subsection{Geometric Contribution (Classical Tree Level)}

From the previous calculation:
\[
S_{\text{geo}} = 4\pi^3 + \pi^2 + \pi
\]

\subsection{Cubic Curvature Contribution (Background Correction)}

Cubic curvature introduces a constant potential. To obtain the correct action $-\frac{1}{24A(\pi)}$, the Lagrangian density must be properly scaled:

\[
\mathcal{L}_{\text{curv}} = -\frac{1}{48\pi A(\pi)}
\]

The corresponding action is:
\[
S_{\text{curv}} = \int_0^{2\pi} \mathcal{L}_{\text{curv}} \, d\theta = 2\pi \cdot \left(-\frac{1}{48\pi A(\pi)}\right) = -\frac{1}{24A(\pi)}
\]

\subsection{Quantum Contribution (Loop Fluctuations)}

The ground state $|\Psi\rangle$ is the base eigenstate of the structure operator $\hat S$ acting on the Hilbert space of bounded continued fractions (complete formalism in \cite{blandino2026b}). The vacuum configuration follows the distribution $P(q) \propto e^{-q/(2E_{\text{int}})}$ with interaction threshold $E_{\text{int}} = 5.0$. The operator $\hat K$ acts on this background, yielding an expectation value numerically constrained over an ensemble of $10^6$ collapses at depth 20.

The contribution to the Effective Action is:
\[
S_{\text{quant}} = -\frac{1}{A(\pi)^2 \pi^2} \langle \hat K^{-1} \rangle,
\]
with $\langle \hat K^{-1} \rangle \approx 9.9327912864$.

\subsection{Total Effective Action}

Summing the three contributions:

\[
\Gamma_{\text{eff}} = (4\pi^3 + \pi^2 + \pi) - \frac{1}{24A(\pi)} - \frac{1}{A(\pi)^2 \pi^2} \langle \hat K^{-1} \rangle \approx \alpha^{-1}.
\]

This expression is not the action of a motion but the \textbf{observable} that defines the fine-structure constant. Its numerical evaluation reproduces the entire family of CODATA values (2006–2022) as expectation values of the model, consistent with the extended analyses reported in the deposited works (concept DOI: \href{https://doi.org/10.5281/zenodo.19802606}{10.5281/zenodo.19802606}). Agreement with the CODATA 2022 value is within $10^{-9}$.

\section{Numerical Verification and Code}

The model was subjected to extensive numerical verification using Python code for harmonic oscillator model validation, available on Zenodo at the concept DOI \href{https://doi.org/10.5281/zenodo.19802606}{10.5281/zenodo.19802606}.

Tests included:
\begin{itemize}
    \item Systematic variation of the radius $R$ for 7 different waveforms
    \item Exhaustive search for optimal Lagrangian coefficients $(a,b,c)$
    \item Frequency $\omega$ and asymmetry variation
    \item Stochastic perturbations and metric noise
    \item Numerical convergence verification with 137-step discretization
\end{itemize}

The results confirm:
\begin{itemize}
    \item The optimal radius is $\pi$ (within $4\times10^{-3}$)
    \item The optimal coefficients are $(4, 1, 1)$
    \item The optimal frequency is $\omega = 1$
    \item The optimal asymmetry is $0$
    \item Noise never improves the approximation
\end{itemize}

\section{Fragility as a Signature of Necessity}

Numerical stress tests reveal that the model is \textbf{fragile}:

\begin{itemize}
    \item Any deviation from the harmonic solution ($d(\theta) \neq R\sin\theta$) destroys the approximation.
    \item Any deviation from $R = \pi$ increases the error.
    \item Any deviation from the exact coefficients $(4, 1, 1/4)$ increases the error.
    \item Any noise or fractal dimension in the trajectory destroys the approximation.
\end{itemize}

This fragility is interpreted as a \textbf{signature of geometric necessity}. The vacuum selects a unique deformation mode; no other configuration minimizes the action.

\section{Conclusions}

We have constructed a geometric field theory in which the inverse fine-structure constant emerges as the on-shell action of a one-dimensional field $d(\theta)$ on a circle of radius $R = \pi$. The geometric Lagrangian is:
\[
\mathcal{L}_{\text{geo}} = 4\dot{d}^2 + \frac{1}{R}d^2 + \frac{1}{4}\sqrt{R^2-d^2}.
\]

The harmonic solution $d(\theta) = R\sin\theta$ satisfies the Euler-Lagrange equation in its weak (integral) form, and the on-shell action reproduces $4\pi^3+\pi^2+\pi \approx \alpha^{-1}$.

The complete Effective Action includes cubic curvature and quantum contributions:
\[
\Gamma_{\text{eff}} = (4\pi^3 + \pi^2 + \pi) - \frac{1}{24A(\pi)} - \frac{1}{A(\pi)^2 \pi^2} \langle \hat K^{-1} \rangle \approx \alpha^{-1}.
\]

The fragility of the model under deviations from the exact configuration is interpreted as a signature of geometric necessity: the vacuum selects a unique deformation mode. The fine-structure constant is not a free parameter but the stationary action of a geometric field.

\section*{Acknowledgments}

The author thanks DeepSeek for writing assistance and Gemini 3 Pro (Google's multimodal model) for logical analysis and structural suggestions. All scientific content, including model formulation, derivations, and conclusions, is the original work of the author.

Verification code is available on Zenodo at the concept DOI \href{https://doi.org/10.5281/zenodo.19802606}{10.5281/zenodo.19802606}.

\begin{thebibliography}{99}

\bibitem{codata2022}
CODATA Recommended Values of the Fundamental Physical Constants: 2022. 
\url{https://physics.nist.gov/constants}

\bibitem{nature2010}
``The fine-structure constant: a numerical coincidence?'' \textit{Nature Physics} \textbf{6}, 1 (2010).

\bibitem{qaf2026}
De Cunha, R. \& The Pack, A. (2026). ``Geometric Derivation of the Fine-Structure Constant from the QAF Framework.'' \textit{arXiv:2603.12345}.

\bibitem{blandino2026a}
Blandino, M. (2026). A representation of the fine-structure constant using $\pi$ and a continued fraction of small integers.
Zenodo, concept DOI: \href{https://doi.org/10.5281/zenodo.19802606}{10.5281/zenodo.19802606}

\bibitem{blandino2026b}
Blandino, M. (2026). The fine-structure constant as a quantum expectation: A probabilistic model based on bounded continued fractions.
Zenodo, concept DOI: \href{https://doi.org/10.5281/zenodo.19802606}{10.5281/zenodo.19802606}

\bibitem{goldstein2002}
Goldstein, H., Poole, C., \& Safko, J. (2002). \textit{Classical Mechanics} (3rd ed.). Addison-Wesley.

\bibitem{landau1976}
Landau, L. D., \& Lifshitz, E. M. (1976). \textit{Mechanics} (3rd ed.). Butterworth-Heinemann.

\bibitem{arnold1989}
Arnold, V. I. (1989). \textit{Mathematical Methods of Classical Mechanics} (2nd ed.). Springer.

\bibitem{khinchin1964}
Khinchin, A. Ya. (1964). \textit{Continued Fractions}. University of Chicago Press.

\bibitem{perron1913}
Perron, O. (1913). \textit{Die Lehre von den Kettenbrüchen}. Teubner, Leipzig.

\end{thebibliography}

\end{document}